% !TEX program = xelatex
% =============================================================================
% 中文求职简历 LaTeX 模板
% Author: Kody
% Repository: https://github.com/kody1126/Chinese-resume-template-work
% Compiler: XeLaTeX
%
% 快速修改指南：
% 1. 在“个人信息头部”中替换姓名、联系方式、求职意向和头像。
% 2. 在“教育背景”中替换学校、专业、成绩排名和课程信息。
% 3. 在实习经历、项目经历、校园经历与荣誉技能模块中增删示例内容。
% 4. 如需调整样式，可修改下方的字体、页面边距与 sectioncolor 配置。
% =============================================================================
\documentclass[11pt,a4paper,fontset=none]{ctexart}

% ── 字体设置 ──────────────────────────────────────────────
% 若需在 Overleaf 使用自定义字体，可将已授权的字体文件放入 fonts/ 文件夹，
% 并按 STSong.ttf、STHeiti.ttf、STKaiti.ttf 命名；macOS 本地默认使用系统华文系列。
\IfFontExistsTF{Times New Roman}{\setmainfont{Times New Roman}}{}

\newif\ifuseuploadedcjkfonts
\IfFileExists{fonts/STSong.ttf}{
    \IfFileExists{fonts/STHeiti.ttf}{
        \IfFileExists{fonts/STKaiti.ttf}{\useuploadedcjkfontstrue}{}
    }{}
}{}

\newif\ifusesystemcjkfonts
\IfFontExistsTF{STSong}{
    \IfFontExistsTF{STHeiti}{
        \IfFontExistsTF{STKaiti}{\usesystemcjkfontstrue}{}
    }{}
}{}

\ifuseuploadedcjkfonts
    \setCJKmainfont[
        Path = fonts/,
        BoldFont = STHeiti.ttf,
        ItalicFont = STKaiti.ttf
    ]{STSong.ttf}
    \setCJKsansfont[Path = fonts/]{STHeiti.ttf}
    \setCJKmonofont[Path = fonts/]{STHeiti.ttf}
\else\ifusesystemcjkfonts
    \setCJKmainfont[
        BoldFont = STHeiti,
        ItalicFont = STKaiti
    ]{STSong}
    \setCJKsansfont{STHeiti}
    \setCJKmonofont{STHeiti}
\else
    \setCJKmainfont[
        BoldFont = Noto Sans CJK SC,
        ItalicFont = Noto Serif CJK SC
    ]{Noto Serif CJK SC}
    \setCJKsansfont{Noto Sans CJK SC}
    \setCJKmonofont{Noto Sans CJK SC}
\fi
\fi

% ── 页面布局 ──────────────────────────────────────────────
\usepackage{geometry}
\geometry{a4paper, left=1.5cm, right=1.5cm, top=1.5cm, bottom=1.5cm}

% ── 颜色定义 ──────────────────────────────────────────────
\usepackage{xcolor}
\definecolor{sectioncolor}{RGB}{25,78,121}
\definecolor{linkcolor}{RGB}{24,90,140}
\definecolor{textcolor}{RGB}{0,0,0}

% ── 标题格式 ──────────────────────────────────────────────
\usepackage{titlesec}
\titleformat{\section}
    {\Large\bfseries\color{sectioncolor}}
    {}
    {0em}
    {\uppercase}
    [\vspace{-0.8em}\rule{\textwidth}{1pt}]
\titlespacing{\section}{0pt}{0.1em}{0.1em}

% ── 列表格式 ──────────────────────────────────────────────
\usepackage{enumitem}
\setlist[itemize]{leftmargin=*, nosep, topsep=0.2em, itemsep=0.1em, parsep=0em}

% ── 图片与排版 ────────────────────────────────────────────
\usepackage{graphicx}
\setkeys{Gin}{draft=false}
\usepackage[export]{adjustbox}
\usepackage{calc}

% ── 表格 ─────────────────────────────────────────────────
\usepackage{tabularx}
\usepackage{multirow}

% ── 图标与超链接 ──────────────────────────────────────────
\usepackage{fontawesome5}
\usepackage{hyperref}
\hypersetup{colorlinks=true, linkcolor=linkcolor, urlcolor=linkcolor}

% ── 全局样式 ──────────────────────────────────────────────
\color{textcolor}
% 中文字符间距：在默认字宽基础上增加 0.15em，使视觉字距约为原来的 1.5 倍。
\xeCJKsetup{CJKglue = \hskip 0.15em}
\linespread{1.00}
\pagestyle{empty}

% ── 自定义命令 ────────────────────────────────────────────
\newcommand{\daterange}[1]{{\textbf{#1}}}
\newcommand{\entrytitle}[1]{{\textbf{#1}}}

% ============================================================
\begin{document}
% ============================================================

% ── 个人信息头部 ──────────────────────────────────────────
\noindent
\begin{minipage}[t]{\textwidth}
    \vspace*{-1.8em}
    \begin{center}
        {\Huge\bfseries\color{sectioncolor} 白逸凡} \\[1.0em]
        \begin{tabular}{@{}ll@{\hspace{2.5em}}ll@{}}
            \textbf{出生日期：} & 2003.09.03 & \textbf{电话：} & 173-2791-2445 \\[0.4em]
            \textbf{个人主页：} & example.com & \textbf{邮箱：} & 1445617868@qq.com \\
        \end{tabular}
    \end{center}
\end{minipage}

\vspace{0.3em}

% ── 教育背景 ──────────────────────────────────────────────
\section*{教育背景}
\vspace{-0.3em}

\noindent
\entrytitle{中国科学院大学信息工程研究所 ｜ 网络信息与安全（硕士）}
\hfill \daterange{2025.09 \textasciitilde{} 至今}

\noindent
\entrytitle{南京邮电大学 ｜ 电子科学与技术（本科）}
\hfill \daterange{2021.09 \textasciitilde{} 2025.06}

\par\noindent
$\bullet$ \textbf{GPA：}3.61 / 5.0（专业前 20\%） \hspace{2em} \textbf{主修课程：}数据结构、计算机组成原理、微机原理、半导体物理与器件等。

% ── 实习经历 ──────────────────────────────────────────────
\section*{实习经历}
\vspace{-0.1em}

\noindent
\entrytitle{北京开源芯片研究院 ｜ SoC 项目组嵌入式软件实习生}
\hfill \daterange{2025.02 \textasciitilde{} 2025.07}

\par\noindent
$\bullet$ \textbf{CPU 测试与工具链：}围绕项目组 CPU 核开展功能调研与测试，研究代码压缩方案，并参与编译器工具链调优。

\par\noindent
$\bullet$ \textbf{车规虚拟化：}基于 ARM-R52 与 bao-hypervisor 调研车机 ECU 芯片虚拟化；梳理 RISC-V 移植所需能力，分析不同 vCPU 调度与动态绑定场景下的功能取舍，掌握特权级虚拟化、内存保护与中断注入等软硬件协同设计思路。

% ── 项目经历 ──────────────────────────────────────────────
\section*{项目经历}
\vspace{-0.1em}

\noindent
\entrytitle{“一生一芯”｜ RISC-V 处理器设计项目成员}
\hfill \daterange{2023.09 \textasciitilde{} 2025.09}

\par\noindent
$\bullet$ \textbf{处理器实现：}完成 RV32IM 指令集模拟器，以 Spike 为 reference 验证正确性；使用 Chisel 以指令组织方式实现 CPU，并在 Verilator 上以自研模拟器为 reference 验证 HDL 实现。

\par\noindent
$\bullet$ \textbf{系统软件与验证：}开发处理器测试工具 Difftest，移植并成功启动 RT-Thread；掌握 RISC-V 指令集、处理器设计与 Chisel 数字电路开发流程。

\par\addvspace{0.65\baselineskip}
\noindent
\entrytitle{全国大学生电子设计大赛（电力电子）｜ 软件成员}
\hfill \daterange{2023.07 \textasciitilde{} 2023.08}

\par\noindent
$\bullet$ \textbf{嵌入式开发：}参与三端口 DC-DC、双向无桥 PFC 与双模式逆变器开发；负责 STM32G474 / TI F280039C 的 ADC、定时器与 DMA 驱动。

\par\noindent
$\bullet$ \textbf{项目成果：}DC-DC 效率不低于 99\%，纹波远低于 5mV；PFC THD 远低于 0.2\%，逆变器效率不低于 98\%，获 2023 年电子设计大赛省一等奖。

\par\addvspace{0.65\baselineskip}
\noindent
\entrytitle{全国大学生集成电路创新创业大赛 ｜ FPGA 成员}
\hfill \daterange{2022.03 \textasciitilde{} 2024.06}

\par\noindent
$\bullet$ \textbf{数字设计与验证：}完成 16×16 多端口 SRAM 控制器 Verilog RTL 与 SVA 验证，45nm 工艺下主频 250MHz、总带宽 16Gbps；参与 JPEG-LS 图像压缩 IP 核 C/C++ Model、Verilog IP 与 AXI4 集成，并在 MPSoC 完成 FPGA 回环测试；获 2023 年省级二等奖。

% ── 科研与开源（可选）：有论文、专利或开源贡献时，取消下方区块注释 ──
% \section*{科研与开源}
% \vspace{-0.1em}
% \noindent \entrytitle{论文 / 专利 / 技术报告标题 ｜ 第一作者（或参与作者）}
% \hfill \daterange{20XX.XX}
% \par\noindent $\bullet$ \textbf{发表信息：}会议 / 期刊 / 专利号 / 技术社区；补充研究问题、核心方法与可复现链接。
% \noindent \entrytitle{开源项目名称 ｜ 维护者 / 核心贡献者}
% \hfill \daterange{20XX.XX \textasciitilde{} 至今}
% \par\noindent $\bullet$ \textbf{开源贡献：}GitHub 链接、Stars / Forks、核心模块、Issue / PR 贡献及实际使用场景。

% ── 校园经历 ──────────────────────────────────────────────
\section*{校园经历}
\vspace{-0.1em}

\noindent
\entrytitle{南京邮电大学大学生科协电子部 ｜ 部长 / 讲师}
\hfill \daterange{2022.09 \textasciitilde{} 2023.06}

\par\noindent
$\bullet$ \textbf{组织与教学：}负责社团日常运营及面向大一学生的授课；策划校“科技节”中的“鲁班电子”与“电子之声”比赛，保障赛程顺利推进。

% ── 荣誉技能：替换证书、技能和获奖经历 ──────────────────
\section*{专业素养}
\vspace{-0.1em}
\noindent $\bullet$ \textbf{技术能力：}C、Python、Scala（Chisel）、Verilog；Linux、单片机开发、PCB 板绘制；Makefile、Verilator、Git、WSL2、Keil C52、示波器。

\noindent $\bullet$ \textbf{荣誉奖项：}二等奖学金、优秀学生干部；EDA 精英挑战赛国家一等奖、电子设计大赛省一等奖、全国集成电路创新创业大赛省二等奖。

\end{document}
